\begin{examplebox}
	\textbf{\textcolor{CExample}{例1（基础）}}

	\textbf{\textcolor{CExample}{题目：}}
	已知函数 $f(x)=x^2-1$，画出 $y=|f(x)|$ 的图像，并指出值域。

	\textbf{\textcolor{CExample}{解题步骤：}}
	\begin{description}[style=nextline, leftmargin=2.8em, labelsep=0.5em, itemsep=0.45em]
		\item[\textbf{第一步（画原图像）：}] 画出 $y=x^2-1$ 的图像，顶点 $(0,-1)$，与 $x$ 轴交于 $(-1,0)$ 和 $(1,0)$，开口向上。
			\par\textbf{理由：} 原二次函数图像是基础。
		\item[\textbf{第二步（识别翻折区域）：}] 图像在 $x$ 轴下方的部分对应 $x\in(-1,1)$，$f(x)<0$；其余部分 $f(x)\ge0$。
			\par\textbf{理由：} 通过解 $x^2-1<0$ 确定区间。
		\item[\textbf{第三步（实施翻折）：}] 将 $x$ 轴下方（$x\in(-1,1)$）的部分关于 $x$ 轴对称翻折到上方，得到 $y=|x^2-1|$ 的图像，形如“W”状。
			\par\textbf{理由：} 应用翻折变换规则。
		\item[\textbf{第四步（写出值域）：}] 图像最低点纵坐标为 $0$（在 $x=\pm1$ 处），向上无限延伸，故值域为 $[0,+\infty)$。
			\par\textbf{理由：} 翻折后图像全部位于 $x$ 轴上方，最小值由原与 $x$ 轴交点处得到。
	\end{description}

	\textbf{\textcolor{CExample}{关键结论：}}
	\textbf{$y=|x^2-1|$ 的值域为 $[0,+\infty)$，图像关于 $y$ 轴对称（偶函数）。}

	\medskip
	\textbf{\textcolor{CExample}{例2（稍变形）}}

	\textbf{\textcolor{CExample}{题目：}}
	求函数 $f(x)=|x^2-4x+3|$ 的单调递增区间。

	\textbf{\textcolor{CExample}{解题步骤：}}
	\begin{description}[style=nextline, leftmargin=2.8em, labelsep=0.5em, itemsep=0.45em]
		\item[\textbf{第一步（分解因式）：}] $x^2-4x+3=(x-1)(x-3)$，所以与 $x$ 轴交于 $(1,0)$ 和 $(3,0)$，开口向上，对称轴 $x=2$。
			\par\textbf{理由：} 因式分解便于找根。
		\item[\textbf{第二步（画出原图像）：}] 画出 $y=x^2-4x+3$ 的图像，顶点 $(2,-1)$，$x\in(1,3)$ 时图像在 $x$ 轴下方。
			\par\textbf{理由：} 基础图像。
		\item[\textbf{第三步（翻折得绝对值图像）：}] 将 $x\in(1,3)$ 的部分关于 $x$ 轴对称翻折，得到 $y=|x^2-4x+3|$ 的图像。
			\par\textbf{理由：} 应用翻折规则。
		\item[\textbf{第四步（观察单调性）：}] 根据翻折后图像，当 $x\le1$ 时，原函数非负且递减，故 $|f(x)|$ 也递减；当 $1\le x\le2$ 时，原函数从 $0$ 递减到 $-1$，绝对值后从 $0$ 递增到 $1$，故递增；当 $2\le x\le3$ 时，原函数从 $-1$ 递增到 $0$，绝对值后从 $1$ 递减到 $0$，故递减；当 $x\ge3$ 时，原函数非负且递增，故 $|f(x)|$ 递增。因此单调递增区间为 $[1,2]$ 和 $[3,+\infty)$。
			\par\textbf{理由：} 利用原函数符号及单调性，结合翻折后函数值变化分析。
	\end{description}

	\textbf{\textcolor{CExample}{关键结论：}}
	\textbf{单调递增区间为 $[1,2]$ 和 $[3,+\infty)$。}

	\medskip
	\textbf{\textcolor{CExample}{例3（综合应用）}}

	\textbf{\textcolor{CExample}{题目：}}
	已知函数 $f(x)=|x^2-2x-3|$ 与直线 $y=k$ 有三个不同的交点，求 $k$ 的取值范围。

	\textbf{\textcolor{CExample}{解题步骤：}}
	\begin{description}[style=nextline, leftmargin=2.8em, labelsep=0.5em, itemsep=0.45em]
		\item[\textbf{第一步（画出绝对值图像）：}] $x^2-2x-3=(x-3)(x+1)$，与 $x$ 轴交于 $(-1,0)$ 和 $(3,0)$，顶点 $(1,-4)$。翻折后得到“W”形图像，最低点为 $(-1,0)$ 和 $(3,0)$，最高点为 $(1,4)$。
			\par\textbf{理由：} 翻折变换应用。
		\item[\textbf{第二步（确定图像关键点）：}] 图像在 $(-\infty,-1]$ 和 $[3,+\infty)$ 上为原二次函数（正值），在 $[-1,3]$ 上为翻折后的凸起，顶点 $(1,4)$，两端 $(-1,0)$ 和 $(3,0)$。
			\par\textbf{理由：} 分析符号及对称性。
		\item[\textbf{第三步（交点个数分析）：}] 作水平直线 $y=k$，根据图像：若 $k<0$，无交点；若 $k=0$，交于 $(-1,0)$ 和 $(3,0)$ 两点；若 $0<k<4$，交于四点（左右各一，中间两个）；若 $k=4$，交于三点（顶点及左右各一）；若 $k>4$，交于两点（左右各一）。因此三个交点对应 $k=4$。
			\par\textbf{理由：} 数形结合，根据图像位置直接判断。
	\end{description}

	\textbf{\textcolor{CExample}{关键结论：}}
	\textbf{$k=4$ 时恰有三个交点。}
\end{examplebox}
